\section{The Harness}
\label{sec:harness}

%% figures/harness.tex -- Figure 1, the harness. Solid boxes are built and
%% measured. Dashed ones belong to the formal track, whose checking core is built
%% but which no formal language or agent has run yet.
\begin{figure}[t]
\centering
\begin{tikzpicture}[
    font=\footnotesize,
    box/.style={draw, rounded corners=2pt, align=center, inner sep=2pt,
                minimum height=10mm},
    planned/.style={box, dashed},
    flow/.style={-{Latex[length=1.5mm]}, semithick},
    back/.style={-{Latex[length=1.5mm]}, densely dotted},
    lbl/.style={font=\scriptsize, inner sep=1pt},
  ]
  \node[box, text width=13mm] (agent) at (0,0) {Agent};
  \node[box, text width=19mm] (gate) at (3.0,0) {PBT gate:\\record, replay,\\4 oracle arms};
  \node[box, text width=14mm] (obj) at (6.0,0) {Objective:\\speed or\\\texttt{val\_bpb}};
  \node[box, text width=76mm, minimum height=7mm] (skills) at (3.0,1.35)
    {\textbf{Skills}: module contracts, review standard, blinded-authoring
     protocol, loop program};
  \node[planned, text width=19mm] (formal) at (3.0,-1.7)
    {\textbf{Formal track}:\\spec, model, proof;\\six checks};
  \node[box, text width=14mm] (ledger) at (6.0,-1.7) {Ledger:\\keep or\\revert};
  \coordinate (bottom) at ($(formal.south) + (0,-2.5mm)$);

  % main flow
  \draw[flow] ([yshift=1.5mm]agent.east) -- ([yshift=1.5mm]gate.west);
  \draw[flow] (gate) -- (obj);
  \draw[flow] (obj) -- (ledger);
  \draw[flow] (skills.south -| agent.north) -- (agent.north);
  \draw[flow, dashed] (gate) -- node[lbl, right] {accepted} (formal);
  \draw[flow, dashed] (formal) -- node[lbl, above] {certificate} (ledger);
  % feedback to the agent
  \draw[back] ([yshift=-1.5mm]gate.west) -- ([yshift=-1.5mm]agent.east);
  \draw[back] (ledger.south) -- (ledger.south |- bottom)
    -- node[lbl, below, pos=0.85] {best so far} ([xshift=-3mm]agent.south |- bottom)
    -- ([xshift=-3mm]agent.south);
  \draw[back, dashed] (formal.west) -| node[lbl, above, pos=0.3] {failed checks}
    ([xshift=3mm]agent.south);
\end{tikzpicture}
\caption{The harness. An agent proposes a candidate with its skills in context. The
  PBT gate scores it on recorded executions, and admitted candidates are measured on the
  objective and kept or reverted. Dotted arrows are feedback to the agent:
  counterexamples from the gate, and the best candidate so far. The formal track
  (dashed) issues certificates only for kernels the gate accepted. The agent writes
  the spec, the model and the proof, and six checks run on PBT's executions decide
  whether the proof counts. Its checking core is built, but no language has run yet.}
\label{fig:harness}
\end{figure}


In each iteration (Figure~\ref{fig:harness}) the gate scores a candidate, and only
admitted candidates are measured on the objective: kernel speed, or validation
bits-per-byte (\texttt{val\_bpb}) when the kernels sit inside a model being trained.

\subsection{Skills}
\label{sec:skills}

A skill is a versioned document an agent loads into its context: a procedure, a set of
rules, or a record of what not to rediscover \citep{anthropic2025skills}. Long sessions
and fresh sub-agents lose what earlier ones measured unless it is written where the next
agent reads it. Four skills carry the weight here.
\textbf{Module contracts}: more than \contractCount{} recorded contracts, each of which
once cost a review round. For example, a reduction's tolerance takes its length
explicitly, and PyTorch's caching allocator hides out-of-bounds accesses from
\texttt{compute-\allowbreak sanitizer}.
\textbf{The review standard}: every test assertion must be the \emph{unique} catcher of
at least one deliberately broken implementation (a saboteur). This is verified by
deleting each assertion in turn.
\textbf{The blinded-authoring protocol}: an agent writing a mutant receives the correct
kernel and one verbatim fault-class description \citep{rathnasuriya2026tile}, and never
the properties, tolerances or results. This skill specifies what is \emph{withheld}.
\textbf{The loop program}: the optimizing agent's instructions name the two cheapest
reward hacks with their measured speedups and state each kernel's semantics.

\subsection{Property-based testing over replayed executions}
\label{sec:pbt}

\paragraph{Record, then replay.}
A seeded generator emits \emph{case groups}: a base input plus metamorphic partners
that share a group identity, such as $x$ and $x + c\mathbf{1}^\top$ with one constant
per row. The whole batch executes once on one backend, and inputs, outputs and device
telemetry are persisted (tensors to safetensors, rows to Parquet). Every oracle is then
evaluated \emph{offline} over the same bytes. This makes the comparison fair by
construction, since no oracle influences which inputs another sees, and a recorded hardware run
can be re-scored by a new oracle without touching a device. Oracles return
\textsc{pass}, \textsc{fail} or \textsc{inconclusive}, so that a case where no
judgment is defined is never counted as a caught bug. The unit of every rate is the
case group.

\paragraph{Four oracle arms.}
The first arm is the field default: \allclose compares with the reference at
$\mathrm{rtol}=\allcloseRtol$ and $\mathrm{atol}=\allcloseAtol$, deliberately
untuned. \RefArm replaces elementwise comparison with a LAPACK-style normalized test ratio
\citep{anderson1992lawn41}:
\begin{equation}
  r(\hat y, y) \;=\; \frac{\lVert \hat y - y \rVert_\infty}
       {\lVert y \rVert_\infty \cdot \eps_{\text{dtype}} \cdot \max(\log_2 n,\, 1)}
  \;\le\; \refThresh,
  \label{eq:ratio}
\end{equation}
where $n$ is the reduction length and $\log_2 n$ is the error bound for pairwise
summation \citep{higham2002accuracy}. \DeclArm evaluates a per-task bundle of laws and never
consults a reference. For softmax the bundle is: every output is finite, every output
lies in $[0,1]$, every row sums to one, and the output is invariant under per-row input
shifts. For layernorm it is: finite outputs, zero row mean, unit row variance and shift
invariance. Tolerances for the bundles come from each dtype's rounding budget. \HybArm
evaluates the laws first and recomputes the reference only when the laws find nothing.

\paragraph{Corpus and backends.}
Development uses a ladder of \ladderGroups{} shapes per task, from $(8,8)$ to $(7,129)$,
with the degenerate rungs $(1,1)$ and $(17,1)$. On single-column rows softmax is
identically $1$ for \emph{any} implementation, so \ladderCeiling{} is the most any arm
can detect on softmax, and every absolute rate below is deflated by that constant. Shift
partners are scaled to $\tfrac12\log(\text{float max})$, the regime where a missing
max-subtraction actually overflows. Kernels run on a NumPy backend and on Triton
\tritonVersion{} on an \gpuName{} (\gpuArch) behind one backend interface.

\subsection{Formal proof: agent-written certificates}
\label{sec:formal}

Proofs serve two purposes here. \emph{Proofs about the oracle} show that a bundle of laws
characterizes its function. They are cheap, needed once per task, and fix what the loop
can converge to (\S\ref{sec:incomplete}). \emph{Proofs about a kernel} show that one
candidate adheres to its spec. They are costly and needed per candidate, and they turn
``admitted'' into ``certified''. The rest of this section is about the second kind. The
kernel-writing agent writes everything
in a formal language: a \emph{spec} formalizing the informal one, a \emph{model} of the
kernel, \emph{hypotheses} under which it claims adherence, and a \emph{proof} of a
theorem whose shape the harness fixes, $\forall x.\; H(x) \Rightarrow S(x, M(x))$. An
agent rewarded for an accepted proof has cheaper routes than proving adherence. It can
weaken the spec, model a simpler kernel, assume a hypothesis nothing satisfies, or
rely on a false axiom. The gate's blind spot (\S\ref{sec:incomplete}) is the first of
these, committed by a person. We therefore let the agent write the spec on purpose and
\emph{measure} the gaming.

%% Source: docs/superpowers/specs/2026-10-01-agent-formal-adherence-design.md, section 5.
\begin{table}[t]
\centering
\small
\setlength{\tabcolsep}{3pt}
\begin{tabular}{@{}lp{0.40\columnwidth}p{0.40\columnwidth}@{}}
\toprule
& Establishes & Sole-catcher cheat \\
\midrule
C1 & checker accepts the proof, no escape hatch & proof ends in \texttt{sorry} \\
C2 & spec accepts the reference's outputs & spec is the kernel, bit for bit \\
C3 & spec rejects every wrong output held, the kernel's own included & spec made of the gate's four laws \\
C4 & model matches the real kernel's executions & model of a tail-dropping kernel computes the whole row \\
C5 & hypotheses hold on every case & $\texttt{BLOCK} \ge n$ on a launch with $\texttt{BLOCK} < n$ \\
C6 & axioms hold on samples; no derivation of \textsf{False} & axiom $\exp(x) = 1 + x$ \\
\bottomrule
\end{tabular}
\caption{The six checks a proof must pass to certify adherence. Each is the only check
  that rejects its cheat. C2--C6 run each artifact's executable reading on PBT's
  cases, with comparisons and tolerances the harness owns.}
\label{tab:checks}
\end{table}

A proof certifies adherence only if all \formalChecks{} checks in Table~\ref{tab:checks}
pass. Every artifact needs an executable reading derived from the proved definitions
themselves, for example a definition over an abstract numeric type instantiated at IEEE
floats. PBT then runs those readings on its own cases, which the agent never sees. C2 and
C3 pin the spec between the reference and a hidden set of wrong kernels that includes
every row of Table~\ref{tab:gap}. Certificates are issued only for kernels the gate
accepted. They never filter the loop, because a proof gate costs speed
\citep{costa2026aicompiler}.

\paragraph{Status.} The core that runs C2--C6 is built and tested against a fake
language adapter. Each check rejects its own cheat and no other, and deleting any of
\formalRemovedConditions{} conditions inside the checks fails a test. Review found three ways to
certify a wrong kernel, each now closed and pinned by a cheat only its fix catches. C3 probed
only the hidden set, so the spec ``$y = \softmax(x)$ or $y = \softmax(1.5x)$'' certified a
$\softmax(1.5x)$ kernel. C4 budgeted at the kernel's own dtype, so a kernel returning
float16 set its own tolerance. And a spec could rewrite the reference or the comparison it
was handed. No formal language or agent has run yet.
